\documentclass[12pt,letterpaper]{article}
\usepackage[utf8]{inputenc}
\usepackage[spanish]{babel}
\usepackage{amsmath}
\usepackage{amsfonts}
\usepackage{amssymb}
\usepackage[margin=2.5cm]{geometry}

\begin{document}

\begin{center}
    \Large \textbf{Evaluación de Examen 2: Unidad II} \\
    \large Física para Ciencias de la Salud (005-1713) \\
    \vspace{0.5cm}
    \textbf{Estudiante:} Valeria Marrero \quad \textbf{C.I.:} 34.264.392
    \vspace{0.3cm} \\
    \large \textbf{Calificación Final: 10/10 puntos}
\end{center}

\vspace{0.5cm}

\section*{Análisis de Resultados}

\subsection*{1. Cinemática (Aceleración media)}
\textbf{Procedimiento y resultado:} Extrae de manera estructurada los datos iniciales y plantea la ecuación de aceleración media ($a = \frac{V_f - V_i}{T}$). Efectúa la sustitución y resta correctas, dividiendo $0,6 \text{ m/s}$ entre $0,15 \text{ s}$. El cálculo aritmético arroja el resultado analítico exacto de $4 \text{ m/s}^2$. \\
\textbf{Veredicto:} \textbf{Correcto} (2/2 pts).

\subsection*{2. Dinámica (Leyes de Newton)}
\textbf{Procedimiento y resultado:} Plantea correctamente la fórmula vectorial de la Segunda Ley de Newton ($\Sigma \vec{F} = m \cdot \vec{a}$) y realiza el despeje impecable para hallar la aceleración ($a = \frac{F}{m}$). Sustituye los valores dividiendo $150 \text{ N}$ entre $25 \text{ kg}$, logrando el resultado perfecto de $6 \text{ m/s}^2$. \\
\textbf{Veredicto:} \textbf{Correcto} (2/2 pts).

\subsection*{3. Trabajo Mecánico}
\textbf{Procedimiento y resultado:} Extrae los datos e identifica la fórmula general del trabajo, incluyendo de forma explícita el componente trigonométrico $\cos(0^\circ)$. Efectúa correctamente la multiplicación ($400 \text{ N} \cdot 0,8 \text{ m}$). El cálculo aritmético da exactamente $320 \text{ J}$. \\
\textbf{Veredicto:} \textbf{Correcto} (2/2 pts).

\subsection*{4. Energía Potencial}
\textbf{Procedimiento y resultado:} Extrae rigurosamente los datos del problema. Plantea el producto de las tres variables según la ecuación de la energía potencial gravitatoria ($E_p = m \cdot g \cdot h$). Destaca muy positivamente que desarrolla el cálculo en pasos secuenciales: halla primero el peso en Newtons ($11,76 \text{ N}$) y luego lo multiplica por la altura ($1,5 \text{ m}$), obteniendo el resultado verificado de $17,64 \text{ J}$. \\
\textbf{Veredicto:} \textbf{Correcto} (2/2 pts).

\subsection*{5. Potencia Mecánica}
\textbf{Procedimiento y resultado:} Extrae los datos y relaciona correctamente el trabajo efectuado y el tiempo planteando la fracción respectiva ($P = \frac{W}{T}$). El desarrollo de la división de $75 \text{ J} / 60 \text{ s}$ arroja la cifra correcta y verificada de $1,25 \text{ W}$. \\
\textbf{Veredicto:} \textbf{Correcto} (2/2 pts).

\vspace{0.5cm}
\section*{Comentarios Generales y Sugerencias}
¡Felicidades por obtener una calificación perfecta! El examen demuestra un nivel de excelencia y madurez analítica sumamente notable.
\begin{itemize}
    \item El formato implementado en cada problema (separando de forma impecable la zona de \textit{Datos} y el desarrollo paso a paso de las ecuaciones a la derecha) facilita enormemente la corrección y el seguimiento metodológico.
    \item Es digno de elogio el estupendo manejo del **análisis dimensional**; la derivación de unidades intermedias como los Newtons en el problema 4 es el estándar de oro en las ciencias de la salud.
    \item Como un pequeñísimo apunte formal de cara a futuros documentos técnicos: la convención estricta del Sistema Internacional indica que el símbolo de kilos (como en kilogramos) se escribe siempre en minúscula ($\text{kg}$ en lugar de $\text{Kg}$).
    \item ¡Sigue con este mismo entusiasmo y nivel de excelencia en las próximas unidades!
\end{itemize}

\end{document}